\Askhsh{1730}{4}
\begin{ylh}{1}
    \item 3.2. 1ο Κριτήριο ισότητας τριγώνων
    \item 4.2. Τέμνουσα δύο ευθειών - Ευκλείδειο αίτημα
    \item 5.2. Παραλληλόγραμμα.
\end{ylh}
Έστω ότι $\mathrm{E}$ και $\mathrm{Z}$ είναι τα μέσα των πλευρών $\mathrm{A}\mathrm{B}$ και $\Gamma\Delta$ παραλληλογράμμου $\mathrm{A}\mathrm{B}\Gamma\Delta$ αντίστοιχα. Αν για το παραλληλόγραμμο $\mathrm{A}\mathrm{B}\Gamma\Delta$ επιπλέον ισχύει $\mathrm{A}\mathrm{B} > \mathrm{A}\Delta$, να εξετάσετε αν είναι αληθείς ή όχι οι ακόλουθοι ισχυρισμοί:\\[0.2cm]
Ισχυρισμός 1: Το τετράπλευρο $\Delta\mathrm{E}\mathrm{B}\mathrm{Z}$ είναι παραλληλόγραμμο. \\
Ισχυρισμός 2: $\mathrm{A}\hat{\mathrm{E}}\Delta = \mathrm{B}\hat{\mathrm{Z}}\Gamma$. \\
Ισχυρισμός 3: Οι $\Delta\mathrm{E}$ και $\mathrm{B}\mathrm{Z}$ είναι διχοτόμοι των απέναντι γωνιών $\hat{\Delta}$ και $\hat{\mathrm{B}}$.
\begin{alist}
    \item Στην περίπτωση που θεωρείτε ότι κάποιος ισχυρισμός είναι αληθής να τον αποδείξετε. \monades{16}
    \item Στην περίπτωση που κάποιος ισχυρισμός δεν είναι αληθής, να βρείτε τη σχέση των διαδοχικών πλευρών του παραλληλογράμμου ώστε να είναι αληθής. Να αιτιολογήσετε την απάντησή σας. \monades{9}
\end{alist}
